\section{Introduction}

Lamb waves are guided elastic waves supported by plates whose two traction-free
surfaces couple longitudinal and shear motion. Their symmetric and antisymmetric
branches are strongly dispersive, so phase velocity, energy transport, and modal
shape all vary with frequency--thickness. These properties make Lamb waves useful
for nondestructive evaluation, while also making broadband measurements difficult
to interpret.

Laser ultrasonics provides non-contact generation and detection. In the
non-destructive thermoelastic regime, absorbed optical energy produces a rapid
temperature rise; thermal expansion then acts as a distributed elastic source.
The temporal laser waveform selects acoustic frequency, its spatial profile
selects wavenumber, and the overlap of the resulting stress with a plate
eigenmode sets the coupling strength.

Most calculations assume a time-harmonic source that has existed indefinitely.
An experiment instead switches on a continuous modulation or applies a finite
burst. This distinction motivates the central question of this work:
\begin{equation}
  \text{When, and in what sense, does the measured field approach its
  frequency-domain prediction?}
\end{equation}
Arrival at $x/\vg$ is not generally equivalent to steady state. The residual
transient samples the curvature of the dispersion surface and becomes especially
important near cutoffs and zero-group-velocity (ZGV) points.

We develop a modal causal formulation and use it to define measurable quantities:
the settling time $t_{\mathrm{ss},n}(x,\epsilon)$, the number of cycles required
to attain a chosen fraction of the continuous-wave amplitude, and the modal
purity during buildup. The same spectral picture naturally extends to translated
and periodically patterned sources.
